% 优化前沿应用：LLM 训练与对齐的横向流程图
% 预训练（SGD/Adam）→ SFT → RLHF（PPO + KL 约束）→ DPO（替代）→ 对齐效果
\documentclass[tikz,border=4pt]{standalone}
\usepackage{ctex}
\usepackage{amsmath,amssymb}
\usetikzlibrary{calc, arrows.meta, positioning, shapes.geometric,
                decorations.pathreplacing, fit, backgrounds}

\begin{document}
\begin{tikzpicture}[scale=1.0, thick,
  stage/.style={draw=#1!70!black, fill=#1!12, rounded corners=5pt,
                minimum width=2.6cm, minimum height=1.4cm,
                font=\footnotesize, align=center, text width=2.4cm},
  arrow/.style={-{Stealth[length=6pt]}, very thick, #1},
  label/.style={font=\tiny, gray!60, align=center},
  detail/.style={draw=gray!30, fill=gray!5, rounded corners=2pt,
                 font=\tiny, align=center, text width=2.2cm},
]

% ─────────────────────────────────────────────────────
% 五个阶段节点（横向排列）
% ─────────────────────────────────────────────────────

% 阶段 1：预训练
\node[stage=gray] (PT) at (0, 3.2)
  {\textbf{预训练}\\
   大规模语料\\
   Adam / Lion};

% 阶段 2：SFT 监督微调
\node[stage=green] (SFT) at (3.2, 3.2)
  {\textbf{SFT 微调}\\
   人工标注对话\\
   AdamW + LoRA};

% RLHF 分支（上方）
\node[stage=orange] (RM) at (6.4, 5.0)
  {\textbf{奖励模型}\\
   Bradley-Terry\\
   偏好对训练};

\node[stage=red] (PPO) at (9.6, 5.0)
  {\textbf{PPO 强化学习}\\
   KL 约束策略\\
   $r(x,y) - \beta\,\text{KL}$};

% DPO 直接路径（下方）
\node[stage=blue] (DPO) at (9.6, 1.4)
  {\textbf{DPO 直接偏好}\\
   无奖励模型\\
   监督学习损失};

% 阶段 5：对齐模型
\node[stage=purple] (ALN) at (12.8, 3.2)
  {\textbf{对齐模型}\\
   有帮助 / 无害\\
   符合人类价值};

% ─────────────────────────────────────────────────────
% 主流程箭头
% ─────────────────────────────────────────────────────

% PT → SFT
\draw[arrow=green] (PT.east) -- (SFT.west)
  node[midway, above, font=\tiny, green!60!black] {规模化语言能力};

% SFT → RM（RLHF 路径）
\draw[arrow=orange] (SFT.north east) to[bend left=15] (RM.west)
  node[midway, above, font=\tiny, orange!70!black, sloped] {RLHF 路径};

% RM → PPO
\draw[arrow=red] (RM.east) -- (PPO.west)
  node[midway, above, font=\tiny, red!70] {$r_\phi(x,y)$ 训练完成};

% PPO → ALN
\draw[arrow=purple] (PPO.south east) to[bend left=15] (ALN.north)
  node[midway, right, font=\tiny, purple!70, sloped] {RL 微调};

% SFT → DPO（直接路径）
\draw[arrow=blue] (SFT.south east) to[bend right=15] (DPO.west)
  node[midway, below, font=\tiny, blue!70, sloped] {DPO 直接路径\\（无 RL loop）};

% DPO → ALN
\draw[arrow=purple] (DPO.east) to[bend right=15] (ALN.south)
  node[midway, right, font=\tiny, purple!70, sloped] {监督微调};

% ─────────────────────────────────────────────────────
% 各阶段详细标注（节点下方）
% ─────────────────────────────────────────────────────

\node[detail] at (0, 1.5)
  {目标：$p_\theta(x)$\\
   Lion / AdamW\\
   $10^{22}\sim10^{24}$ FLOPs};

\node[detail] at (3.2, 1.5)
  {目标：学对话格式\\
   LoRA / 全参微调\\
   $10^{20}$ FLOPs};

\node[detail] at (6.4, 6.8)
  {$r_\phi(x,y_w) > r_\phi(x,y_l)$\\
   BTL 模型训练};

\node[detail] at (9.6, 6.8)
  {$\max\mathbb{E}[r] - \beta\text{KL}$\\
   PPO clip；KL 惩罚};

\node[detail] at (9.6, -0.1)
  {$-\log\sigma[\beta(\Delta\log\pi)]$\\
   2 个模型，无 RL};

\node[detail] at (12.8, 1.5)
  {有帮助 (Helpful)\\
   无害 (Harmless)\\
   诚实 (Honest)};

% ─────────────────────────────────────────────────────
% 路径分叉标注
% ─────────────────────────────────────────────────────
\node[draw=gray!40, fill=white, font=\tiny, align=center, rounded corners]
  at (6.4, 3.2)
  {分叉点\\RLHF $\uparrow$\\DPO $\downarrow$};

% ─────────────────────────────────────────────────────
% 优化器汇总（顶部）
% ─────────────────────────────────────────────────────
\node[draw=gray!30, fill=gray!5, rounded corners=3pt,
      font=\tiny, align=left, text width=5.5cm]
  at (0.5, 7.2)
  {\textbf{关键优化器}\\
   预训练：Adam / Lion（显存$\downarrow$1/3）/ Sophia（步数$\downarrow$2$\times$）\\
   微调：AdamW + LoRA（显存$\downarrow$10$\sim$100$\times$）\\
   对齐：AdamW（标准）；DPO 直接监督学习};

% ─────────────────────────────────────────────────────
% 标题
% ─────────────────────────────────────────────────────
\node[font=\small\bfseries] at (6.4, 8.1)
  {大语言模型训练与对齐的优化前沿流程};

\end{tikzpicture}
\end{document}
